\chapter{罪数}

% Source PDF physical page 171; see audit/source_coverage.json.
\begin{EditionOpeningPanel}
理论上将罪数分为“本来的一罪”“科刑的一罪”等。法考不考查这些分类标准。因此，本书仅从答题角度予以讲解。做题第一步应先判断行为的数量，是一个行为还是两个行为。
\end{EditionOpeningPanel}

\section{一个行为}

[问题]甲在夜市弹吉他卖唱，强迫顾客听歌，不给钱，就殴打，顾客被迫给钱。甲构成强迫交易罪还是抢劫罪？（卖唱案）

\subsection*{一、继续犯}

继续犯也称为持续犯，是指犯罪行为与不法状态在一定时间内一直处于持续状态的犯罪，只定一罪。例如，非法拘禁罪、窝藏罪、持有型犯罪（非法持有枪支罪、非法持有毒品罪等）。

[注意］继续犯在犯罪既遂后，行为终了前，第三人参与犯罪的，成立承继的共同犯罪。这一点与一般犯罪不同。一般犯罪既遂后，有人参与，只可能构成事后的妨害司法犯罪（窝藏、包庇罪，掩饰、隐瞒犯罪所得罪，帮助毁灭证据罪）。

例1，甲盗窃到财物，构成盗窃罪既遂。此后乙帮甲保管财物。乙构成掩饰、隐瞒犯罪所得罪，而不构成盗窃罪的帮助犯。

例2，甲将丙非法拘禁后，构成非法拘禁罪既遂。此后，乙参与进来一起非法拘禁丙。乙构成非法拘禁罪的承继的共同犯罪。

\subsection*{二、法条竞合}

例如，刑法既规定了诈骗罪，又规定了贷款诈骗罪。当甲诈骗银行的贷款时，既符合诈骗罪的条文，又符合贷款诈骗罪的条文。诈骗罪与贷款诈骗罪属于法条竞合的关系。

1.法条关系：一般与特别的关系

一般与特别的关系，也即A与A+B的关系，B是特殊因子。这种关系也称为包容关系，也即A+B可以包容评价为A。例如，诈骗罪是A（他人财产权），贷款诈骗罪是A+B（B是金融秩序）。贷款诈骗罪不仅侵犯了他人财产权，还破坏了金融秩序。贷款诈骗罪可以包容评价为诈骗罪，触犯贷款诈骗罪就必然触犯诈骗罪。

% Source PDF physical page 172; see audit/source_coverage.json.
2.判断标准：触犯的必然性

\CJKunderwave{若触犯此罪名，必然会触犯彼罪名，则二者是法条竞合关系}。例如，触犯合同诈骗罪，必然会触犯诈骗罪，因此二者是法条竞合关系。前者是特殊罪名，后者是一般罪名。又如，触犯盗窃罪并不必然触犯故意杀人罪，因此二者不是法条竞合关系。

3.处理规则：特殊法优于一般法。例如，甲的一个行为同时触犯贷款诈骗罪和诈骗罪，要优先适用贷款诈骗罪。

\subsection*{三、想象竞合犯}

这是指，一个行为同时触犯两个罪名。处理规则是：择一重罪论处。例如（偷救命药案），甲带着非法占有目的，盗窃乙的急救药，知道这样做可能会导致乙死亡，仍然盗窃，乙无药可救而死亡。甲的一个行为同时触犯盗窃罪和故意杀人罪，两罪想象竞合，择一重罪论处，最终定故意杀人罪。

1.法条竞合与想象竞合的关系

相同点：都是一个行为，触犯两个罪名。

区别：两个罪名的关系不同。法条竞合中，两个罪名的关系是A与A+B的包容关系。想象竞合中，两个罪名的关系是A与B的中立关系（不搭噶的关系）。

\CJKunderwave{具体判断标准：触犯的必然性不同}。触犯此罪名就必然触犯彼罪名，则两个罪名属于法条竞合关系。如果没有这种必然性，才有可能是想象竞合。例如，触犯盗窃罪并不必然触犯故意杀人罪，触犯故意杀人罪并不必然触犯盗窃罪，因此二者不是法条竞合关系，要构成想象竞合关系，需要想象出来一个行为，将二者同时触犯，例如上文“偷救命药案”。\footnote{法条竞合中，两个罪名是永恒的同时触犯。想象竞合中，两个罪名是临时的同时触犯。例如，诈骗罪（A)与合同诈骗罪（A+B）是永恒的同时触犯。又如，盗窃罪（A）与故意杀人罪（B），二者在正常情况下是不搭噶的，只是由于某个犯罪人实施了某个行为，临时将二者同时触犯（如偷救命药案）。

法条竞合，是立法规定的特殊产物。例如，立法者在诈骗罪之外又规定了合同诈骗罪，由此产生法条竞合现象。想象竞合，不是立法规定的特殊产物，而是罪名之间的正常现象，是原生态（原本状态）的现象。因此，如果按照法条竞合来处理会出现问题时，可以退回原生态（原本状态），按照想象竞合来处理。}

2.法条竞合与想象竞合的动态变化

两个法条竞合关系的罪名，按照特殊法优于一般法来处理，如果会导致罪刑不相适应，此时可以将法条竞合关系调整为想象竞合关系，按照择一重罪论处。这种特殊情况主要出现在下面情形中：一般罪名与特殊罪名的犯罪形态不一致。

例如，一般情况下，触犯故意杀人罪必然触犯故意伤害罪，两罪是法条竞合关系，故意杀人罪是特殊罪名，故意伤害罪是一般罪名，优先以故意杀人罪论处。但是，甲杀害乙，导致乙重伤时，甲又中止。甲同时构成故意杀人罪中止和故意伤害罪（重伤）既遂。如果一定要以故意杀人罪中止论处，有可能判得比故意伤害罪（重伤）既遂轻，这显然不合适。这时候，只能认为故意杀人罪中止与故意伤害罪（重伤）既遂是想象竞合关系，择一重罪论处。

% Source PDF physical page 173; see audit/source_coverage.json.
3.罪名关系的类型

\InsertDiagram{figures/manual_table_p173_1.tex}

4.罪名关系的注意点

（1）不要将A与B的关系当作A与-A的关系，也即不要将中立关系当作对立关系。

例如，关于本节开头的“卖唱案”中，传统理论认为，强迫交易罪与抢劫罪的区分标准是，“是否发生在市场交易背景下，如是，则定强迫交易罪；如不是，则定抢劫罪”。如此，将两罪视为A与一A的关系，将两罪对立排斥起来。由于有卖唱这个行业背景，因此对甲定强迫交易罪。这显然是不妥当的。这等于鼓励抢劫犯通过卖唱去抢劫，如此能规避抢劫罪。正确看法是，两罪是A与B的中立关系，可以想象竞合。甲的一个行为同时触犯强迫交易罪和抢劫罪，想象竞合，择一重罪论处，定抢劫罪。

（2）不要将A与B的关系当作A与A+B的关系，也即不要将中立关系当作包容评价关系。

例如，甲使用普通诈骗手法骗取小芳一千万元。乙冒充官员，骗取小芳一千万元。甲构成诈骗罪，可判处无期徒刑。对于乙，传统理论认为，乙同时构成诈骗罪和招摇撞骗罪，二者属于法条竞合，招摇撞骗罪是特殊罪名。这种观点的问题是，第一，法条竞合要求“触犯特殊罪名，必然触犯一般罪名”，例如，触犯合同诈骗罪，必然触犯诈骗罪。而触犯招摇撞骗罪，并不必然触犯诈骗罪。例如，狗蛋冒充官员招摇撞骗，只骗色，不骗财。狗蛋触犯招摇撞骗罪，但未触犯诈骗罪。第二，法条竞合的处理规则是特殊法优于一般法。基于此，应优先认定为特殊罪名，也即定招摇撞骗罪。而对招摇撞骗罪，最高只能判十年有期徒刑。甲乙都骗了一千万元，乙冒充官员诈骗，反倒判得很轻。\footnote{传统理论为此又“打补丁”，认为对这种法条竞合可以特殊处理，也即按照择一重罪论处。然而，“择一重罪论处”是想象竞合的处理规则。因此，不如大大方方承认，两罪是想象竞合关系。}实际上，诈骗罪的保护法益是财产，招摇撞骗罪保护的法益是国家机关工作人员的公众信赖感。两罪的保护法益不同，是A与B的关系。因此，两罪不是法条竞合关系，可以想象竞合。乙同时构成诈骗罪和招摇撞骗罪，想象竞合，择一重罪论处，以诈骗罪论处，可判处无期徒刑。

（3）不要将A与A+B的区分关系当作A与-A的区分关系，也即不要将包容评价的关系当作对立排斥关系。

% Source PDF physical page 174; see audit/source_coverage.json.
例如（2017年第60题的类似案例），乙站在街头抽烟，身旁停一辆摩托车，车主是丙，丙不在现场。甲以为乙是车主，想抢劫这辆车，上前将乙打倒，迅速骑走该车。甲主观上想犯抢劫罪，客观上是盗窃罪。盗窃罪与抢劫罪是性质相同、程度不同的关系，也即A与A+B的包容评价关系。抢劫罪可以包容评价为盗窃罪，在盗窃罪的范围内，主客观相一致，对甲定盗窃罪。传统理论认为，抢劫罪必须是公开的，盗窃罪必须是秘密的。由于“公开”与“秘密”是反义词，是A与一A排斥关系，那么公开的抢劫罪便无法包容评价为秘密的盗窃罪，对甲便无法定盗窃罪。这显然是不合理的。

\subsection*{四、加重犯}

加重犯是本讲最疑难的知识点，需要勇敢的你耐心学习，多看两遍。

\subsubsection*{（一）行为结构}

加重犯是一个行为，同时触犯两个罪名，刑法将一个罪名作为基本犯，将另一个罪名作为基本犯的法定刑升格条件，予以加重处罚。例如，抢劫罪是基本犯，“抢劫罪致人死亡”是结果加重犯，“入户抢劫”是情节加重犯，“抢劫数额特别巨大的财物”是数额加重犯，“抢劫军用物资”是对象加重犯。这些加重犯都是抢劫罪的法定刑升格条件。

1.加重犯是一个行为，不是两个行为。例如，情节加重犯“入户抢劫”。从自然事实角度看，甲实施了两个行为，“非法侵入住宅”和“抢劫”。但是，在法律评价上，\CJKunderwave{一个}\CJKunderwave{目的之下实施的两个举动，都被评价为一个行为}。甲带着非法取财目的，实施了非法侵入住宅和劫财行为，在整体上实施了一个抢劫行为，“非法侵入住宅”属于抢劫行为的一部分。

2.加重犯是法律的特殊规定，其原生态是想象竞合。例如，抢劫罪是A，非法侵入住宅罪是B，入户抢劫是A+B。甲带着非法取财目的入户抢劫。甲的一个行为同时触犯A罪与B罪，二者原本是想象竞合关系，择一重罪论处，但是，法律规定成情节加重犯“A+B”。

3.正常情况下，基本犯与加重犯是法条竞合关系。例如，抢劫罪是A，非法侵入住宅罪是B，入户抢劫是A+B。触犯“A+B”罪必然触犯A罪，所以二者是法条竞合关系，也即情节加重犯“入户抢劫”与基本犯“抢劫罪”是法条竞合关系。按照特殊法优于一般法的规则，以“入户抢劫”论处。

注意：这里有两组不同关系，一是A罪与B罪的关系，二者原本是想象竞合关系，但法律特殊规定为情节加重犯“A+B”。二是A罪与A+B罪的关系，二者是法条竞合关系。\footnote{注意，虽然B罪与A+B罪也是法条竞合关系，但是不能认为：“非法侵入住宅罪（B罪）是基本犯，入户抢劫（A+B）是非法侵入住宅罪的情节加重犯”。这是因为，“入户抢劫”没有规定在非法侵入住宅罪里，没有规定为非法侵入住宅罪的法定刑升格条件。}

4.表述重心问题。加重犯包括两部分，一是基本犯的部分（基础部分），二是加重部分。注意，当称加重犯时，评价重心是加重部分。例如，当称“入户抢劫”是情节加重犯，评价重心是“入户”，也即带着抢劫目的而非法入户，而非“抢劫”部分，“抢劫”部分属于基本犯的内容。又如，当称“强奸罪致人重伤”是结果加重犯，评价重心是

% Source PDF physical page 175; see audit/source_coverage.json.
“致人重伤”，也即带着奸淫目的实施暴力致人重伤，而非“强奸”部分，“强奸”部分属于基本犯的内容。

\subsubsection*{（二）基本犯与加重犯的犯罪形态与关系调整}

正常情况下，基本犯（A罪）与加重犯（A+B）是法条竞合关系。这里所谓的正常情况，是指基本犯与加重犯的犯罪形态是一致的。

基本犯都是故意犯罪，存在未遂、中止、既遂形态。当加重犯的加重部分是故意行为，便存在未遂、中止、既遂形态。当基本犯的犯罪形态与加重犯的加重部分的犯罪形态不一致，此时，基本犯与加重犯的法条竞合关系就要调整为想象竞合关系。\footnote{实际上，可以将加重犯当作一个独立罪名（加重型罪名，加重构成要件）。例如，盗窃枪支罪可称为盗窃罪的对象加重犯，只是立法者没有将其规定为盗窃罪的法定刑升格条件，而规定成独立罪名。甲在军营盗窃枪支，将箱子偷回家，发现是普通财物。甲的一个行为，一方面构成盗窃罪（基本犯，既遂），一方面构成盗窃枪支罪（加重犯，未遂），二者是想象竞合关系，择一重罪论处。无论立法者将盗窃枪支规定为独立罪名，还是对象加重犯，都不改变“想象竞合择一重”这种处理结论。}这是为了确保遵守罪刑相适应原则。

1.基本犯既遂，加重犯未遂或中止。

例1（情节加重犯），甲夜晚在公园对乙女实施暴力，欲当众强奸乙，乙哀求不要在公园强奸，甲便将乙拉到附近灌木丛强奸，奸淫了乙。基本犯强奸罪构成既遂。甲也构成情节加重犯“在公共场所当众强奸妇女”，属于情节加重犯的中止。二者想象竞合，择一重罚方案论处。有可能基本犯的既遂处罚更重，有可能加重犯的中止处罚更重，哪个处罚更重，就定哪个。

[结论］当基本犯既遂，加重犯未遂或中止，那么想象竞合，择一重罚方案论处。

例2（数额加重犯），甲欲盗窃某珠宝店3万元的珠宝，手伸进玻璃柜（此时已经属于盗窃着手），临时认错对象，拿走旁边价值3千元的珠宝。甲的一个行为，一方面构成基本犯“盗窃数额较大”的既遂，一方面构成数额加重犯“盗窃数额巨大”的未遂，想象竞合，择一重罚方案论处。

2.基本犯未遂或中止，加重犯既遂。

例1（结果加重犯），甲为了强奸妇女，刚将妇女打成重伤，就被赶到的警察抓捕。甲成立“强奸罪（故意）致人重伤”，属于结果加重犯的加重部分（故意致人重伤）的既遂。但基本犯强奸罪是未遂。二者想象竞合，择一重罚方案论处。由于结果加重犯的既遂的处罚更重，所以定结果加重犯的既遂，也即适用“十年至死刑”的升格法定刑，再适用未遂犯的规定，也即可以从宽处罚。\footnote{之所以再适用未遂犯的规定，是因为升格法定刑是在基本犯既遂的基础上确立的、升格的。当基本犯是未遂、中止，那么在适用升格法定刑时，应再适用未遂犯、中止犯的规定，予以从宽处罚。}

[结论］当基本犯未遂或中止，加重犯既遂，那么定加重犯既遂，适用升格法定刑，再适用未遂犯、中止犯的规定，从宽处罚。

例2（情节加重犯），甲为了抢劫，非法入户，刚实施暴力，就被抓。甲成立“入户抢劫”，属于情节加重犯的加重部分（带着抢劫目的而非法入户）既遂。但基本犯抢劫罪是未遂。定情节加重犯的既遂，也即适用“十年至死刑”的升格法定刑，再适用未遂犯的规定，也即可以从宽处罚。

% Source PDF physical page 176; see audit/source_coverage.json.
\section{两个行为}

[问题]甲盗窃到一辆车，来到二手车市场，谎称是自已合法的车，用正常行情价卖给乙，乙不知情而买下。对甲如何处理？（卖赃车案）

\subsection*{一、结合犯}

1.概念：两个原本独立的犯罪，根据法律规定结合成一个犯罪。例如，狗蛋为拐卖小芳，将小芳控制到手，又强奸了小芳，然后卖掉。按照正常做法，狗蛋构成拐卖妇女罪和强奸罪，数罪并罚。但是，刑法特殊规定：只定拐卖妇女罪，将强奸罪作为拐卖妇女罪的法定刑升格条件。强奸罪被拐卖妇女罪吞并了。

2.结合形式：甲罪+乙罪=甲罪，乙罪成为甲罪的法定刑升格条件，乙罪被甲罪吞并了。结合犯必须由法律明文规定。这是结合犯的法定性。如果没有法定性，对两罪应并罚。

3.结合犯与加重犯的关系

二者的相同点：都是法律规定的特殊产物，都是法定刑升格条件。例如，“奸淫被拐卖的妇女”是拐卖妇女罪的法定刑升格条件。“拐卖致人重伤、死亡”也是拐卖妇女罪的法定刑升格条件。

二者的区别：（1）加重犯是一个行为，不是两个行为；若删除法律规定，则还原的原生态是想象竞合犯。（2）结合犯是独立的两个行为；若删除法律规定，则还原的原生态是数罪并罚。

例1，狗蛋拐卖小芳，为了制服小芳，殴打小芳，导致小芳重伤。狗蛋构成拐卖妇女罪致人重伤（结果加重犯）。如果删除这个结果加重犯的规定，则狗蛋的一个殴打行为，一方面构成拐卖妇女罪，另一方面构成故意伤害罪，想象竞合，择一重罪论处。

例2，狗蛋使用暴力将小芳拐到手，在出卖途中，强奸了小芳。狗蛋构成“奸淫被拐卖的妇女”（结合犯）。如果删除这个结合犯的规定，则狗蛋的两个行为，暴力行为构成拐卖妇女罪，强奸行为构成强奸罪，数罪并罚。

[总结]我国刑法中的结合犯：

（1）拐卖罪+强奸罪=拐卖罪，加重处罚。

（2）绑架罪+故意杀人罪=绑架罪，加重处罚；绑架罪+故意伤害罪致人重伤、死亡=绑架罪，加重处罚。（2008年第8题）

（3）前罪+妨害公务罪=前罪，加重处罚。这样的前罪只有：组织、运送他人偷越国（边）境罪，走私、贩卖、运输、制造毒品罪。其他罪+妨害公务罪，均应数罪并罚。（2006年第7题）

（4）交通肇事罪+“因逃逸致人死亡”（遗弃罪致人死亡）=交通肇事罪，加重处罚。

% Source PDF physical page 177; see audit/source_coverage.json.
\subsection*{二、连续犯}

这是指，基于同一种犯意，连续实施性质相同的数个行为，触犯同一个罪名。\footnote{这里的连续犯是指广义的连续犯，包括了同种数罪。狭义的连续犯与同种数罪的区别在于，是否在较短时间内实施。例1，甲在某小区内，同一天内连续诈骗三户人家，属于狭义的连续犯。例2，甲在一月份诈骗一户人家，在二月份诈骗一户人家，在三月份诈骗一户人家，不属于狭义的连续犯，而属于同种数罪，属于广义的连续犯。}对连续犯的处理规则：

1.如果罪名的保护法益不具有人身专属性，则定一罪。

例1，甲连续多次实施诈骗罪。诈骗罪的保护法益是财产，财产法益不具有人身专属性，具有可替代性，因此定一个诈骗罪，数额累计计算。

例2，甲多次走私，多次逃税，多次贪污，多次受贿。这些罪的保护法益不具有人身专属性，定一个罪，数额累计计算。

例3（行为对象的转移），甲欲盗窃乙家的现金，发现又有珠宝，便放弃现金，拿走珠宝。甲对现金构成盗窃罪中止，对珠宝构成盗窃罪既遂。但是，不需要并罚，只需要定一个盗窃罪既遂即可，因为财产法益不具有人身专属性，而具有可替代性。（2010年第57题）

2.如果罪名的保护法益具有人身专属性，多数说认为，原则上应数罪并罚。\footnote{参见张明楷：《论同种数罪的并罚》，载《法学》2011年第1期，第127页。}

例1，甲进入某个办公楼，每一层杀一个人，连续杀了三个人，属于连续犯，应数罪并罚。因为，故意杀人罪的保护法益是生命法益，具有人身专属性，具有不可替代性。如果甲在多地连环杀人，在一年内杀了三个人，更应数罪并罚。\footnote{数罪并罚的一个优点是，能够妥当处理自首问题。如果只定一个故意杀人罪，问题在于，甲如果自动投案，但只如实供述了第一个故意杀人罪，隐瞒了后两个故意杀人罪，对甲是否认定自首，便存在障碍。若认定自首，但甲没有如实供述后两个故意杀人罪。若不认定自首，但甲如实供述了第一个故意杀人罪。如果数罪并罚，则不存在这个障碍，也即每个故意杀人罪是否构成自首，各算各的账，对第一个故意杀人罪认定为自首，对后两个故意杀人罪不认定为自首。}

例2（行为对象的转移），甲本欲强奸女生乙，进入宿舍按倒乙后，发现同寝室的丙更漂亮，便放过乙，强奸了丙。甲对乙构成强奸罪中止，对丙构成强奸罪既遂，应并罚，因为人身法益具有专属性。

3.如果法条规定“多次实施”按照一罪处理，则按照一罪处理。

\InsertDiagram{figures/manual_table_p177_1.tex}

\subsection*{三、吸收犯}

这是指，实施了两个行为，重行为吸收了轻行为，仅成立重行为的罪名。例如，甲盗窃枪支后非法持有，非法持有行为被盗窃行为吸收，只定盗窃枪支罪。

% Source PDF physical page 178; see audit/source_coverage.json.
吸收犯不是法律规定的产物。\CJKunderwave{重行为吸收轻行为的条件}：不遗漏评价法益侵害。如果做不到这一点，就必须数罪并罚。

例1，甲过失致乙重伤，又故意不救乙，导致乙死亡。前行为是过失致人重伤罪，后行为是不作为的故意杀人罪。故意杀人罪可以吸收过失致人重伤罪。

例2，甲盗掘古墓葬时，将乙的炕挖塌，乙掉进坑里。甲不救乙，乙死亡。前行为是盗掘古墓葬罪，后行为是不作为的故意杀人罪。由于法益性质不同，因此故意杀人罪不能吸收盗掘古墓葬罪，应数罪并罚。

[巩固练习]题1，甲制造枪支后又持有枪支，对甲只需定非法制造枪支罪。乙制造毒品后又持有毒品，对乙只需定制造毒品罪。（2010年第55题）

题2，甲伪造假币后又出售。甲先构成伪造货币罪，后构成出售假币罪，最终只定重罪（伪造货币罪）。

题3，甲先敲诈乙，乙不怕甲，甲便用暴力打倒乙，然后取财。甲先构成敲诈勒索罪未遂，后构成抢劫罪既遂。由于两个行为具有先后的紧密联系，且都是财产犯罪，所以重罪吸收轻罪。

题4，甲从超市偷到一部手机，又要求“退货”，换取到现金。甲先构成盗窃罪，后构成诈骗罪。由于甲最终获得了一份好处，超市遭受了一份损失，所以没必要并罚，重罪吸收轻罪。一般而言，盗窃罪更重，所以最终定盗窃罪。

题5，甲打算入户抢劫，进入户内，发现没人，拿走财物。甲在预备阶段构成抢劫罪（犯罪预备），在实行阶段构成盗窃罪既遂。实行行为的处罚一般重于预备行为，实行行为吸收预备行为，最终定盗窃罪既遂。

\subsection*{四、不可罚的事后行为}

这是指某个犯罪已经既遂，又实施了另一个犯罪行为，但是不处罚事后行为。不处罚的根据是：第二个行为不具有期待可能性或者没有侵犯新的法益。不可罚的事后行为实际上是吸收犯的一种情形，因为太重要，故单独论述。

1.不具有期待可能性。例如（2008年第8题），甲盗窃到车辆后，在赃车渠道销售，买家知道是赃车，心照不宣。对这种销售行为不具有期待可能性，因为无法期待小偷偷到财物后不去销赃，所以对这种销赃行为不处罚，也即不定掩饰、隐瞒犯罪所得罪。

2.没有侵犯新的法益

例1（2008年第8题），甲以为盗窃到普通财物，发现是枪支，甲继续非法持有。甲构成盗窃罪和非法持有枪支罪，并罚。

例2（2011年第56题），甲盗窃到一幅名画后，误以为是品便毁坏。因为只侵犯了一个法益（名画的财产所有权），所以不再定故意毁坏财物罪，只定盗窃罪。

[巩固练习]题1，甲盗窃到一幅名画后，发现是品，便冒充真迹欺骗他人，骗取1万元。因为又侵犯了另一人的财产权，所以定盗窃罪和诈骗罪，并罚。（2010年第55题，2019年试题）

% Source PDF physical page 179; see audit/source_coverage.json.
题2，本节开头“卖赃车案”中，甲的盗窃行为构成盗窃罪。甲的销售行为欺骗乙，导致乙有财产损失，也即用正常行情价买到一个赃车，让乙花了冤枉钱，侵犯了乙的法益，构成诈骗罪（多数说）。两罪应并罚。

[总结］盗窃者销赃有两种不同情形：第一种是正常销赃，属于不可罚的事后行为。第二种是欺骗型销赃，对买家构成诈骗罪。

\subsection*{五、牵连犯}

1.概念：两个犯罪行为具有手段与目的的关系，或者原因与结果的关系，这种关系被称为牵连关系。例如，伪造国家机关公文（手段行为），然后用来骗取财物（目的行为）。

2.这种牵连关系要具有通常性的特征。不是任何手段行为与目的行为结合起来就构成牵连犯，只有手段行为与目的行为在实践中经常结合在一起，具有通常性特征，才能构成牵连犯。

例1（2010年第19题），为了诈骗而伪造、购买国家机关证件，具有通常性特征，属于牵连犯，择一重罪论处。

例2（2016年第54题），甲盗窃军人制服，然后冒充军人招摇撞骗，触犯盗窃罪和冒充军人招摇撞骗罪，不具有通常性特征，不是牵连犯，应数罪并罚。

[巩固练习]题1（2017年第8题），甲盗窃国家机关证件，然后持该证件招摇撞骗，触犯盗窃国家机关证件罪和招摇撞骗罪，不具有通常性特征，不是牵连犯，应数罪并罚。

题2，为了实施信用卡诈骗，伪造信用卡，触犯伪造金融票证罪和信用卡诈骗罪，具有通常性特征，属于牵连犯，择一重罪论处。

题3，为了杀人去盗窃枪支或抢劫他人的杀猪刀，不具有通常性特征，不是牵连犯，应数罪并罚。

3.牵连犯与想象竞合犯的区别：前者有两个行为，后者只有一个行为。

4.处理方法：择一重罪论处，并且按照该重罪论处时，还可以从重处罚。

[总结]罪数原理

\InsertDiagram{figures/manual_table_p179_1.tex}

% Source PDF physical page 180; see audit/source_coverage.json.
\ContinuedTable{figures/manual_table_p180_1.tex}

\OriginalPanel{典型真题}

关于罪数的处理，下列说法错误的有？（2019年回忆版金题）\footnote{[答案]A项，触犯“轮奸”必然触犯强奸罪。因此二者是法条竞合关系，基本犯（强奸罪第一款）是一般法条，情节加重犯（轮奸，也即强奸罪第三款）是特别法条。B项，甲的出卖行为构成诈骗罪，侵犯了新的法益，因此不属于不可罚的事后行为，需要与前面的盗窃罪并罚。C项，（1）钱某触犯的第一种情节加重犯是“入户抢劫”。

（2）钱某触犯的第二种情节加重犯是“持枪抢劫”。（3）钱某触犯的第三种情节加重犯是“多次抢劫”。“多次抢劫”是指抢劫3次或以上，并且每次均构成抢劫罪。（4）虽然钱某触犯了三种情节加重犯，但对它们不能并罚，否则意味着重复处罚了，只能择一重罚方案论处。第一种方案是，定一个抢劫罪，然后按照“多次抢劫”加重处罚。不过，在此就不能再以“入户抢劫”“持枪抢劫”处罚，否则会重复处罚了。第二种方案是，定第一个抢劫罪，属于“入户抢劫”；定第二个抢劫罪，属于“入户抢劫”；定第三个抢劫罪，属于“持枪抢劫”；然后三个抢劫罪并罚。不过在此就不能再以“多次抢劫”处罚。两种方案的处罚一样重，都能够到死刑。所以，选择哪个方案都行。C项说法错误。D项，周某不构成抢劫罪致人死亡。周某构成抢劫罪和故意杀人罪，应并罚。本题答案为ABCD。}

A.“二人以上轮奸”只是强奸罪的法定刑升格条件，与强奸罪的关系不是特别法条与一般法条的关系

B.甲发现自己盗窃到的是一件仿真品（价值4000元），冒充真品以2万元卖给他人。甲的变卖行为是不可罚的事后行为

C.钱某分别实施了两次入户抢劫，一次持枪抢劫，均构成抢劫罪。钱某分别触犯了三种抢劫罪的情节加重犯，对这三种情节加重犯应并罚

D.周某抢劫了陈某的财物后，担心暴露，杀害了陈某。周某构成抢劫罪致人死亡和故意杀人罪的想象竞合

